\documentclass[a4paper]{article}

\usepackage{graphicx}
\usepackage{amsmath,amssymb}
\usepackage{minted}
\usepackage{multirow}
\usepackage{float}
\usepackage{todonotes}
\usepackage{forest}
\usepackage{xstring}
\usepackage{xcolor}
\usepackage[most]{tcolorbox}
\usepackage{xparse}
\usepackage{natbib}

\usetikzlibrary{graphs,graphdrawing}
\usegdlibrary{layered}

% for external graphviz
\input{/home/donkarlo/Dropbox/repo/nd_language_ai_project/src/nd_language_ai/formal/computer/programming/tex/latex/code_clips/external_graphviz.tex}

% for tags
\input{/home/donkarlo/Dropbox/repo/nd_language_ai_project/src/nd_language_ai/formal/computer/programming/tex/latex/parts/control_sequence/kind/semtag/semtag.tex}

% for links
\input{/home/donkarlo/Dropbox/repo/nd_language_ai_project/src/nd_language_ai/formal/computer/programming/tex/latex/code_clips/seehere.tex}

\title{Cell signaling}
\date{}
\author{Mohammad Rahmani}

\begin{document}
    \maketitle
    Cell signaling is the biological process by which a cell interacts with itself, with other cells, and with the environment. Cell signaling is a fundamental property of all forms of life. . Typically, the signaling process involves three components: the first messenger (the ligand), the receptor, and the signal itself \seeherefooter{https://en.wikipedia.org/wiki/Cell_signaling}.
    
    \begin{itemize}
        \item Neuron signaling at the end of the axon is an example
        \item Insoline release is another example of cell signalling
    \end{itemize}
    
    \textbf{Other resources}
    \begin{itemize}
        \item Cell to Cell Communication || Types of signaling \seeherefooter{https://www.youtube.com/watch?v=i3bY-JCYs4A}
        \item intracell communication \seeherefooter{https://www.youtube.com/watch?v=aIZQ3ker0KE}
    \end{itemize}
    
    \bibliographystyle{plainnat}
    \bibliography{/home/donkarlo/Dropbox/repo/nd_shared_research/bibliography/bibliography.bib}
\end{document}